\documentclass[10pt,a4paper]{amsart}
\usepackage[margin=19mm]{geometry}
\usepackage{amsmath,amssymb,mathtools}
\usepackage[scr=rsfs,cal=euler]{mathalfa}
\usepackage{xfrac}
\usepackage[unicode,hidelinks,bookmarks=false]{hyperref}
\newcommand{\ie}{i.e.\ }
\newcommand{\cf}{cf.\ }
\newcommand{\resp}{respectively\ }
\newcommand{\Subsection}{\subsection}
\providecommand{\nobreakdash}{\nobreak\hbox{-}}
\setlength{\emergencystretch}{2em}
\multlinegap0pt
\usepackage{bm}
\usepackage{mathtools}
\usepackage{amssymb}
\usepackage[shortlabels]{enumitem}
\usepackage{dsfont}
\usepackage{xcolor}
\usepackage{tikz}
\usepackage{csquotes}
\newtheorem{proposition}{Proposition}
\newtheorem{lemma}{Lemma}
\newtheorem{theorem}{Theorem}
\usepackage{cleveref}
\crefname{proposition}{Proposition}{Propositions}
\crefname{lemma}{Lemma}{Lemmas}
\crefname{theorem}{Theorem}{Theorems}
\newcommand{\M}{\op{M}}
\newcommand{\br}{\mathbb R}
\newcommand{\op}{\operatorname}
\title{Determinant values on lattices}
\author{Wooyeon Kim, Hee Oh}
\date{Source: arXiv:2607.18038v2 (CC BY 4.0)}
\begin{document}
\maketitle

\noindent\emph{Source and licence.} Determinant values on lattices. Version arXiv:2607.18038v2 (CC BY 4.0), distributed by arXiv under the Creative Commons Attribution 4.0 International licence, http://creativecommons.org/licenses/by/4.0/, as recorded in the arXiv licence metadata for this exact version and shown on the abstract page https://arxiv.org/abs/2607.18038v2. Full licence notice: \texttt{licenses/arxiv-2607.18038-CC-BY-4.0.txt}.\par\smallskip

\noindent\textit{Editorial scope.} Counted unit: the article's theorem thm:globalcontraction (source0.tex lines 4903--4920). The article's complete original proof of it is delivered with it (source0.tex lines 4922--5309, one \texttt{proof} environment of 388 lines). No mathematics is written, repaired or re-derived: the statement and the proof are the article's own lines, in the article's own notation, and no retained line is substituted or re-flowed.  Retained uncounted as the source-local context the proof needs: lines 3652--3670; lines 3798--3803; lines 3818--3829; lines 4068--4071; lines 4198--4204; lines 4254--4260; lines 4317--4321; lines 4323--4327; lines 4332--4343; lines 4648--4650; lines 4730--4761.  Nothing was omitted from any retained span, so the package carries no omission record.  No retained line contains a citation, so the wrapper carries no bibliography and no bibliography entry.  No retained line contains a figure environment, an image-inclusion command or drawing code, so no image is delivered and the package declares no asset dependency.  Editorial formatting only, in addition: an AMS \texttt{amsart} preamble that restates the article's own declarations and macros used by the retained lines, namely the environments \texttt{proposition}, \texttt{lemma}, \texttt{theorem} on separate counters, together with the standard packages those lines need.  The retained environments print in this document as Proposition 1, Lemma 1, Theorem 1 in order of appearance, whereas the article prints the same items with the numbering of its own full document.  The complete native source of the article as distributed in the arXiv e-print for this exact version is bundled unchanged as \texttt{sources/205622/source0.tex}.
\begin{proposition}
\label{lem:localcontractioninequalitygenerali}
\label{cor:generalweightedcontractiontensor}
Let $\Omega\subset\mathbb R^{n-1}\times\mathbb R^{n-1}$ be bounded and
open. For every
$\frac12\le\beta\le1+\frac1{2n^2}$, there exists $c>0$ such that
for every $t\ge1$ and nonzero $v\in\mathcal V_{\op{gen}}$, \begin{equation}
\label{eq:contractioninequalityatleasthalf}
\int_\Omega\|a_tu_{\boldsymbol{\xi}} v\|^{-\beta}\,d\boldsymbol{\xi}
\le c e^{-t/(5n^2)}\|v\|^{-\beta}.
\end{equation}
 Moreover, for all
$0\le\beta\le1+\frac1{2n^2}$,
\begin{equation}
\label{eq:contractioninequalitysmallrange}
\int_\Omega\|a_tu_{\boldsymbol{\xi}} v\|^{-\beta}\,d\boldsymbol{\xi}
\ll_{\beta,\Omega}\|v\|^{-\beta}.
\end{equation}
\end{proposition}

\begin{equation}
\label{eq:phi-as-minimum}
\phi_{kn}(v)
=
\min\{F_{00}(v),F_{12}(v),F_{21}(v),F_{10}(v),F_{20}(v)\}.
\end{equation}

\begin{proposition}[Contraction for $\phi_{kn}$]
\label{prop:localMargulisineq;n-1}
Let $1\le k\le n-1$, $0<\theta<(2n)^{-6}$, and
$1\le\beta\le1+\frac{\theta}{4n^2}$.
There exists $c=c(\beta,\theta)>0$ such that, for every
$t\ge1$ and 
$v\notin\mathcal M_{k,1}\cup\mathcal M_{k,2}$,
$$
\int_{B_{U}(1)}\phi_{kn}(a_tuv)^\beta\,du
\le c e^{-\theta t/(20n^2)}\phi_{kn}(v)^\beta.
$$
\end{proposition}

\begin{equation}\label{eq:beta-tau-comparison}
\beta_i^{-1}<\tau_i
\qquad (1\le i\le N-1).
\end{equation}

\begin{equation}\label{eq:sharp-majorized-by-tilde}
\widehat{\alpha}_{i,\eta,M}(h;\Delta)^{\beta_i}
\ll_n \widetilde\alpha_{i,\eta,M}(h;\Delta),
\qquad
\widehat{\alpha}_{\eta,M}(h;\Delta)^{\beta_*}
\ll_n \widetilde\alpha_{\eta,M}(h;\Delta).
\end{equation}

\begin{equation}\label{eq:Youngineq}
x^{\beta_i/2}y^{\beta_i/2}z
\ll
\rho\left(x^{\beta_{i+j}}+y^{1/\beta_{i-j}}\right)
+
1+\left(\rho^{-1}z\right)^{C_n\theta_1^{-1}}.
\end{equation}

\begin{equation}\label{eq:logLipschitzforwedge}
 e^{-4n^3 s}\|v\|
\le \|a_suv\|
\le e^{4n^3 s}\|v\|.
\end{equation}

\begin{equation}\label{eq:logLipschitzformodifiedht}
e^{-4n^3s}\phi_{kn}(v)
\le \phi_{kn}(a_suv)
\le e^{4n^3s}\phi_{kn}(v)
\end{equation}

\begin{lemma}[Margulis inequality for $\bar{\alpha}$]\label{prop:globalboundedexpansion'}
There are constants $c_n,C_n>0$ such that, for every $\Delta\in X$ and
$s\ge 2\log n$
\begin{equation}\label{eq:globalboundedexpansion'}
\int_{B_{U}(1)}
\bar\alpha(a_su\Delta)\,du
\ll
 e^{-c_n\theta_1 s}\bar\alpha(\Delta)
+
 e^{C_n\theta_1^{-1}s}.
\end{equation}
\end{lemma}

\begin{equation}\label{eq:phi-decomposition-global}
\phi_i(v)=\|v\|^{-1+4\theta_i}\bar\phi_i(v)^{4\theta_i}.
\end{equation}

\begin{lemma}\label{lem:BQinequality}
Let $\Delta\in X$, and let $W_1,W_2$ be distinct
$\Delta$-rational subspaces of dimension $i=kn$ such that $W_1+W_2\ne \M_n(\br)$. Let $j$ be such that 
$$
\dim(W_1\cap W_2)=i-j\quad 
\text{ and }\quad
\dim(W_1+W_2)=i+j.
$$
Let $w_m=\mathsf w_{\Delta,W_m}$, $m=1,2$, and assume that
$w_m\notin\mathcal M_{k,1}\cup\mathcal M_{k,2}$ and
$$
\|\pi_{k,1}(w_m)\|
=\max_{q=0,1,2}\|\pi_{k,q}(w_m)\|
\qquad (m=1,2).
$$
Suppose also that, for some $\mathsf L\ge1$,
$$
\mathsf L^{-1}\|w_2\|\le\|w_1\|\le\mathsf L\|w_2\|\;\; \text{ and }\;\; 
\mathsf L^{-1}\bar\phi_i(w_2)
\le\bar\phi_i(w_1)
\le\mathsf L\bar\phi_i(w_2).
$$
Then
\begin{equation}\label{eq:BQ-modified-height-inequality}
\phi_i(w_1)\phi_i(w_2)
\ll
\mathsf L^{8\theta_i}
\|\mathsf w_{\Delta,W_1\cap W_2}\|^{-1}
\phi_{i+j}(\mathsf w_{\Delta,W_1+W_2}).
\end{equation}
The same conclusion holds with $\pi_{k,1}$ replaced by $\pi_{k,2}$.
\end{lemma}

\begin{theorem}[Main global Margulis inequality outside the exceptional set]\label{thm:globalcontraction}
Fix $b\ge1$ and $M'\ge1$. There exists
$\theta_0=\theta_0(n,b,M')>0$ such that for every
$0<\theta\le\theta_0$, $0<\eta<1$, $M\ge1$,
$s\ge 2\log n$ and
$(h,\Delta)\notin\mathcal E_{s,\eta,M,M'}$,
\begin{equation}\label{eq:globalcontraction}
\begin{aligned}
\int_{B_U(1)}
\widetilde\alpha(a_suh;\Delta)\,du
\ll{}&
e^{-c_n\theta s}
\widetilde\alpha(h;\Delta)
\log\!\bigl(3+\widehat\alpha_{\eta,M}(h;\Delta)\bigr)
+e^{C_n(M'\theta+\theta^{-1})s}.
\end{aligned}
\end{equation} Here the implicit constant depends only on $b$, $M'$, and $\theta$.
\end{theorem}

\begin{proof}
By reducing $\theta_0$, we may assume that
\begin{equation}\label{eq:theta-small-global-contraction}
C_nbM'\theta\le\frac1{20}.
\end{equation}

All local integral estimates below are applied on a fixed bounded open
neighborhood of $B_{U}(1)$ and then restricted to $B_{U}(1)$.
Fix $1\le i\le N-1$. Let
$\mathcal P_i=\mathcal P_i(h,\Delta,s)$ be the finite collection
of $i$-dimensional $\Delta$-rational subspaces $V$ such that
\begin{equation}\label{eq:relevant-subspaces-definition}
V\notin\widetilde{\mathscr Q}_{i,\eta,M}(\Delta),
\quad
0<\|h\mathsf w_{\Delta,V}\|\le1,
\quad 
\phi_i(h\mathsf w_{\Delta,V})^{\beta_i}
>e^{-12n^3s}\widetilde\alpha_{i,\eta,M}(h;\Delta).
\end{equation}
Finiteness follows because the corresponding primitive vectors in
$\wedge^i(h\Delta)$ lie in a bounded set.
We first bound the contribution of subspaces outside
$\mathcal P_i$. Let $W$ be any $\Delta$-rational subspace occurring
in the supremum in
$\widetilde\alpha_{i,\eta,M}(a_suh;\Delta)$, and write $w:=\mathsf w_{\Delta,W}$. Thus $W\notin\widetilde{\mathscr Q}_{i,\eta,M}(\Delta)$ and $0<\|a_suhw\|\le1$, and its contribution to the supremum is
$\phi_i(a_suhw)^{\beta_i}$. By
\eqref{eq:logLipschitzforwedge}, we have $\|hw\|\le e^{4n^3s}$.

If
$\|hw\|\le1$ and $W\notin\mathcal P_i$, then
\eqref{eq:logLipschitzformodifiedht} and
\eqref{eq:relevant-subspaces-definition} give
\begin{equation}
\begin{aligned}
        \phi_i(a_suhw)^{\beta_i}
&\le e^{4n^3\beta_i s}\phi_i(hw)^{\beta_i}
\\&\le e^{(4n^3\beta_i-12n^3)s}
\widetilde\alpha_{i,\eta,M}(h;\Delta)
\le e^{-4n^3s}\widetilde\alpha_{i,\eta,M}(h;\Delta),
\end{aligned}
\end{equation}
because $\beta_i<2$.

If $\|hw\|>1$ and $n\nmid i$, the same expression is
$O(e^{2C_ns})$. Finally, suppose $i=kn$ and $\|hw\|>1$. Since
$W\notin\widetilde{\mathscr Q}_{i,\eta,M}(\Delta)$, its own Pl\"ucker vector does not
belong to $\mathscr Q_{kn,\eta,M}$. Let
\begin{equation}\label{eq:exceptional-scale}
    \varepsilon_{s,\eta,M,M'}(h;\Delta)
:=
\begin{cases}
e^{-M's}
&\text{if } \alpha(h\Delta)
\le e^{20n^3s},\\ 
\alpha(h\Delta)^{-bM'}
&\text{if }\alpha(h\Delta)
>e^{20n^3s}.\end{cases}
\end{equation}  Because
$(h,\Delta)\notin\mathcal E_{s,\eta,M,M'}$,
$$
\min_{m=1,2}\|hw-\pi_{k,m}(hw)\|
>\varepsilon_{s,\eta,M,M'}(h;\Delta).
$$
Since $\|hw\|>1$ and $-1+4\theta_i<0$, the definition of $\phi_i$ gives
$\phi_i(hw)^{\beta_i}\ll
\varepsilon_{s,\eta,M,M'}(h;\Delta)^{-4\theta_i\beta_i}$.
Using \eqref{eq:logLipschitzformodifiedht} and the two alternatives in
\eqref{eq:exceptional-scale},
we obtain
\begin{equation}\label{eq:outside-P-error}
\phi_i(a_suhw)^{\beta_i}
\ll
e^{2C_ns}
\left(
 e^{C_nM'\theta s}
 +\alpha(h\Delta)^{C_nbM'\theta}
\right).
\end{equation}
With the convention that the maximum over the empty set is zero, it follows that
\begin{equation}\label{eq:competitor-reduction}
\begin{aligned}
\widetilde\alpha_{i,\eta,M}(a_suh;\Delta)
\ll{}&
\max_{V\in\mathcal P_i}
\phi_i(a_suh\mathsf w_{\Delta,V})^{\beta_i}
+e^{-4n^3s}\widetilde\alpha_{i,\eta,M}(h;\Delta)\\
&+e^{2C_ns}
\left(e^{C_nM'\theta s}+\alpha(h\Delta)^{C_nbM'\theta}\right).
\end{aligned}
\end{equation}

Fix a sufficiently large constant $C_0=C_0(n)$. 

\noindent{\bf{Case (1).}} First suppose that
\begin{equation}\label{eq:few-competitors}
\#\mathcal P_i
\le 3C_0\log(3+\widehat{\alpha}_{\eta,M}(h;\Delta)).
\end{equation}
For noncritical $i$, every summand of $\wedge^i\M_n(\mathbb R)$ is
generic, and \Cref{lem:localcontractioninequalitygenerali}
applies. For $i=kn$, use
\Cref{prop:localMargulisineq;n-1}. Thus
$$
\int_{B_{U}(1)}
\phi_i(a_suhv)^{\beta_i}\,du
\ll e^{-c_n\theta_1 s}\phi_i(hv)^{\beta_i}
$$
for every vector $v=\mathsf{w}_{\Delta,V}$ with $V\in \mathcal P_i$. Integrating \eqref{eq:competitor-reduction} and using
\eqref{eq:few-competitors} gives the required decaying term. For the error
term, the definition of $\bar\alpha$ gives $\alpha(h\Delta)
\le \bar\alpha(h\Delta)^{1/\tau_*}
\le \widetilde\alpha_{\eta,M}(h;\Delta)^{1/\tau_*}$.
Together with \eqref{eq:theta-small-global-contraction} and
$\tau_*>1/2$, this implies
$$
\alpha(h\Delta)^{C_nbM'\theta}
\ll 1+\widetilde\alpha_{\eta,M}(h;\Delta)^{1/2}.
$$
The inequality
$x\widetilde\alpha_{\eta,M}(h;\Delta)^{1/2}
\le \rho\widetilde\alpha_{\eta,M}(h;\Delta)+C\rho^{-1}x^2$, with
$\rho=e^{-c_n\theta_1 s}$, absorbs this sublinear power. Consequently,
\begin{multline}\label{eq:few-competitors-bound}
\int_{B_{U}(1)}
\widetilde\alpha_{i,\eta,M}(a_suh;\Delta)\,du
\\ \ll
 e^{-c_n\theta_1 s}\widetilde\alpha_{\eta,M}(h;\Delta)
 \log(3+\widehat{\alpha}_{\eta,M}(h;\Delta))
+e^{C_n(M'\theta+\theta_1^{-1})s}.
\end{multline}

\noindent{\bf{Case (2).}} Now suppose that
\begin{equation}\label{large}
\#\mathcal P_i
> 3C_0\log(3+\widehat{\alpha}_{\eta,M}(h;\Delta)).
\end{equation} If $n\nmid i$, put
$\mathcal R_i^0=\mathcal P_i$ and
$\mathcal R_i^1=\mathcal R_i^2=\emptyset$. If $i=kn$, partition
$\mathcal P_i$ into three sets. Let $\mathcal R_i^0$ consist of
those $V$ for which
$$
\phi_i(h\mathsf w_{\Delta,V})
\le e^{4n^3s}\|h\mathsf w_{\Delta,V}\|^{-1}.
$$
If the generic component of $h\mathsf w_{\Delta,V}$ were maximal, then
\eqref{eq:phi-as-minimum} and equivalence of the fixed norms would give
$\phi_i(h\mathsf w_{\Delta,V})\asymp_n
\|h\mathsf w_{\Delta,V}\|^{-1}$. Since $s\ge2\log n$, such a subspace
belongs to $\mathcal R_i^0$. Thus every remaining subspace has a maximal
exceptional component.
Place each remaining subspace in $\mathcal R_i^1$ or
$\mathcal R_i^2$ according to whether a maximal component of
$h\mathsf w_{\Delta,V}$ is the $\pi_{k,1}$-component or the
$\pi_{k,2}$-component, assigning ties to $\mathcal R_i^1$. At least one of
these three sets contains more than
$C_0\log(3+\widehat{\alpha}_{\eta,M}(h;\Delta))$ elements.


\noindent{\bf{Case (2-a)} $\#\mathcal R_i^0\ge 2$.} 
Let $W_1\ne W_2$ be two subspaces in $\mathcal R_i^0$. Set
$$
\dim(W_1\cap W_2)=i-j,
\qquad
\dim(W_1+W_2)=i+j.
$$
If $i+j<N$ and
$W_1+W_2\in
\widetilde{\mathscr Q}_{i+j,\eta,M}(\Delta)$,
then the downward-closure property and
$W_1\subset W_1+W_2$ would imply
\[
W_1\in\widetilde{\mathscr Q}_{i,\eta,M}(\Delta),
\]
contrary to the definition of $\mathcal P_i$. Therefore
\[
W_1+W_2
\notin
\widetilde{\mathscr Q}_{i+j,\eta,M}(\Delta).
\]
When $i+j=N$, we use the endpoint convention introduced above.
The Pl\"ucker covolume inequality and
\eqref{eq:relevant-subspaces-definition} give
\begin{equation}\label{eq:R0-product-bound}
\widetilde\alpha_{i,\eta,M}(h;\Delta)^{2/\beta_i}
\ll
e^{2C_ns}
\alpha_{i-j}^+(h\Delta)
\widehat\alpha_{i+j,\eta,M}(h;\Delta).
\end{equation}
Taking the $\beta_i/2$-power of
\eqref{eq:R0-product-bound} and absorbing the additional log-Lipschitz
factor from \eqref{eq:competitor-reduction} into $e^{C_ns}$, we obtain
$$
e^{2C_ns}\widetilde\alpha_{i,\eta,M}(h;\Delta)\ll e^{2C_ns}
\widehat\alpha_{i+j,\eta,M}(h;\Delta)^{\beta_i/2}
\alpha_{i-j}^+(h\Delta)^{\beta_i/2}.
$$
Apply \eqref{eq:Youngineq} with
$x=\widehat\alpha_{i+j,\eta,M}(h;\Delta)$,
$y=\alpha_{i-j}^+(h\Delta)$, $z=e^{2C_ns}$, and
$\rho=e^{-c_n\theta_1s}$, omitting an endpoint variable when necessary.
This gives
$$
e^{2C_ns}\widetilde\alpha_{i,\eta,M}(h;\Delta)
\ll
 e^{-c_n\theta_1 s}
\left(
\widehat\alpha_{i+j,\eta,M}(h;\Delta)^{\beta_{i+j}}
+
\alpha_{i-j}^+(h\Delta)^{1/\beta_{i-j}}
\right)
+e^{C_n\theta_1^{-1}s},
$$
with the absent endpoint term omitted. By \eqref{eq:beta-tau-comparison} and
\eqref{eq:sharp-majorized-by-tilde}, both displayed height terms are
$\ll \widetilde\alpha_{\eta,M}(h;\Delta)$. Consequently,
\begin{equation}\label{eq:R0-contraction}
e^{2C_ns}\widetilde\alpha_{i,\eta,M}(h;\Delta)
\ll e^{-c_n\theta_1 s}\widetilde\alpha_{\eta,M}(h;\Delta)
+e^{C_n\theta_1^{-1}s}.
\end{equation}
This settles the case (2-a).

\noindent{\bf{Case (2-b)} $\#\mathcal R_i^0\le1$.} 

Since $\mathcal P_i$ satisfies \eqref{large}, either
$\mathcal R_i^1$ or $\mathcal R_i^2$ then contains more than
$C_0\log(3+\widehat{\alpha}_{\eta,M}(h;\Delta))$ elements. By symmetry, assume that
$\mathcal R_i^1$ is large. Every
$W\in\mathcal R_i^1$ satisfies
$$
\widehat{\alpha}_{\eta,M}(h;\Delta)^{-1}
\le\|h\mathsf w_{\Delta,W}\|\le1.
$$
The interval
$[\widehat{\alpha}_{\eta,M}(h;\Delta)^{-1},1]$ is covered by at most
$\ll \log(3+\widehat{\alpha}_{\eta,M}(h;\Delta))$ dyadic
subintervals. Since $C_0$ was chosen sufficiently
large, the pigeonhole principle gives distinct
$W_1,W_2\in\mathcal R_i^1$ whose Pl\"ucker norms are comparable by an
absolute factor. By \eqref{eq:relevant-subspaces-definition} and the
definition of $\widetilde\alpha_{i,\eta,M}(h;\Delta)$, the
$\phi_i$-value of each lies between
$e^{-12n^3s/\beta_i}\widetilde\alpha_{i,\eta,M}(h;\Delta)^{1/\beta_i}$ and
$\widetilde\alpha_{i,\eta,M}(h;\Delta)^{1/\beta_i}$. They are therefore comparable by a factor
$e^{2C_ns}$. Using \eqref{eq:phi-decomposition-global} and the comparability
of their Pl\"ucker norms, we obtain
\begin{equation}\label{eq:barphi-comparison-competitors}
\mathsf L^{-1}\bar\phi_i(hw_2)
\le\bar\phi_i(hw_1)
\le\mathsf L\bar\phi_i(hw_2) \quad\text{for }\mathsf L:= e^{2C_ns/\theta_i},
\end{equation}
 where $w_m=\mathsf w_{\Delta,W_m}$. Set
$$
w_-:=\mathsf w_{\Delta,W_1\cap W_2},
\qquad
w_+:=\mathsf w_{\Delta,W_1+W_2}.
$$
If $i+j<N$, then, as above,
$W_1+W_2\notin\widetilde{\mathscr Q}_{i+j,\eta,M}(\Delta)$.

\noindent{\bf{Case (2-b-(i))}} Suppose first that $i+j<N$ and $\|hw_+\|\le1$.  Apply
\Cref{lem:BQinequality} to the lattice $h\Delta$ and the subspaces
$hW_1,hW_2$. Together with
\eqref{eq:barphi-comparison-competitors}, this yields
$$
\phi_i(hw_1)\phi_i(hw_2)
\ll e^{2C_ns}
\|hw_-\|^{-1}\phi_{i+j}(hw_+).
$$
Combining this with \eqref{eq:relevant-subspaces-definition} gives
$$
\widetilde\alpha_{i,\eta,M}(h;\Delta)^{2/\beta_i}
\ll
e^{2C_ns}
\alpha_{i-j}^+(h\Delta)
\phi_{i+j}(hw_+).
$$
Because $W_1+W_2\notin\widetilde{\mathscr Q}_{i+j,\eta,M}(\Delta)$ and
$\|hw_+\|\le1$, the last factor satisfies
$$
\phi_{i+j}(hw_+)
\le
\widetilde\alpha_{i+j,\eta,M}(h;\Delta)^{1/\beta_{i+j}}.
$$

Consequently, after absorbing the additional log-Lipschitz factor from
\eqref{eq:competitor-reduction} into $e^{C_ns}$, we obtain
$$
e^{2C_ns}\widetilde\alpha_{i,\eta,M}(h;\Delta)\ll e^{2C_ns}
\widetilde\alpha_{i+j,\eta,M}(h;\Delta)^{\beta_i/(2\beta_{i+j})}
\alpha_{i-j}^+(h\Delta)^{\beta_i/2}.
$$
Apply \eqref{eq:Youngineq} with
$x=\widetilde\alpha_{i+j,\eta,M}(h;\Delta)^{1/\beta_{i+j}}$,
$y=\alpha_{i-j}^+(h\Delta)$, $z=e^{2C_ns}$, and
$\rho=e^{-c_n\theta_1s}$, using its one-factor version when $i-j=0$.
Using
\eqref{eq:beta-tau-comparison}, we obtain
\begin{equation}\label{eq:R1-Mother-contraction}
e^{2C_ns}\widetilde\alpha_{i,\eta,M}(h;\Delta)
\ll e^{-c_n\theta_1 s}\widetilde\alpha_{\eta,M}(h;\Delta)
+e^{C_n\theta_1^{-1}s}.
\end{equation}

\noindent{\bf{Case (2-b-(ii))}} Now suppose that either $\|hw_+\|>1$ or $i+j=N$. In the latter
case, $hw_+$ is a top-degree Pl\"ucker vector of the unimodular lattice
$h\Delta$, and hence $\|hw_+\|=1$. Thus, in either case, the standard intersection-sum inequality and the comparability of $\|hw_1\|$ and
$\|hw_2\|$ imply
$$
\|hw_-\|\ll\|hw_1\|^2.
$$
Hence $\|hw_1\|^{-1}\ll
\alpha_{i-j}^+(h\Delta)^{1/2}$. In the region where the first
exceptional component is maximal, norm equivalence gives
$$
\phi_i(hw_1)\asymp_n
\|hw_1\|^{-1+2\theta_i}
\|hw_1-\pi_{k,1}(hw_1)\|^{-2\theta_i}.
$$
Because $(h,\Delta)$ is outside the exceptional set,
$$
\|hw_1-\pi_{k,1}(hw_1)\|
>\varepsilon_{s,\eta,M,M'}(h;\Delta).
$$
Using the formula for $\phi_i$ in the region where the first exceptional
component is maximal, we obtain
$$
\phi_i(hw_1)
\ll
\alpha_{i-j}^+(h\Delta)^{1/2}
\left(
 e^{2M'\theta_is}
 +\alpha(h\Delta)^{2bM'\theta_i}
\right).
$$
The relevance condition for $W_1$, together with
$\alpha_{i-j}^+(h\Delta)
\le \widetilde\alpha_{\eta,M}(h;\Delta)^{1/\tau_{i-j}}$ when $i-j\ge1$
(and $\alpha_0^+=1$ otherwise) and
$\alpha(h\Delta)\le \widetilde\alpha_{\eta,M}(h;\Delta)^{1/\tau_*}$, now gives
$$
\widetilde\alpha_{i,\eta,M}(h;\Delta)
\ll
e^{2C_ns}e^{C_nM'\theta s}
\widetilde\alpha_{\eta,M}(h;\Delta)^{q_i},
$$
where $q_i:=
\frac{\beta_i}{2\tau_{i-j}}
+
\frac{2bM'\theta_i\beta_i}{\tau_*}$ if $i-j\ge 1$ and
$q_i:=
\frac{2bM'\theta_i\beta_i}{\tau_*}$ if $i-j=0$.

There are only finitely many admissible pairs $(i,j)$. As $\theta\to0$,
$\beta_i$ and $\tau_*$ tend to $1$, as does $\tau_{i-j}$ when $i-j\ge 1$,
while $bM'\theta_i$ tends to $0$.
Thus $q_i$ tends to $1/2$, or to $0$ when $i-j=0$.
After decreasing $\theta_0(n,b,M')$, we therefore have $q_i\le2/3$
uniformly. Hence
\begin{equation}\label{eq:R1-endpoint-sublinear}
\widetilde\alpha_{i,\eta,M}(h;\Delta)
\ll e^{2C_ns}e^{C_nM'\theta s}
\widetilde\alpha_{\eta,M}(h;\Delta)^{2/3}.
\end{equation}
The elementary scaled inequality
$x\widetilde\alpha_{\eta,M}(h;\Delta)^{2/3}
\le \rho\widetilde\alpha_{\eta,M}(h;\Delta)+C\rho^{-2}x^3$, with
$\rho=e^{-c_n\theta_1 s}$, therefore gives the same conclusion as
\eqref{eq:R1-Mother-contraction}. The case of $\mathcal R_i^2$ is
identical.

In Case (2), the log-Lipschitz estimate gives
$$
\max_{V\in\mathcal P_i}
\phi_i(a_suh\mathsf w_{\Delta,V})^{\beta_i}
\ll e^{2C_ns}\widetilde\alpha_{i,\eta,M}(h;\Delta).
$$
Combining this with \eqref{eq:competitor-reduction},
\eqref{eq:R0-contraction}, \eqref{eq:R1-Mother-contraction}, and
\eqref{eq:R1-endpoint-sublinear}, and absorbing the same sublinear error as
in Case (1), yields \eqref{eq:few-competitors-bound}. Thus this
bound holds for every $1\le i\le N-1$. Summing over $i$ and adding
\Cref{prop:globalboundedexpansion'} gives the corresponding estimate in
terms of $\theta_1$. Since $\theta_1=10^{1-N}\theta$ and $N=n^2$, replacing
$\theta_1$ by $\theta$ only changes the dimensional constants. Absorbing
this change into $c_n$ and $C_n$ yields \eqref{eq:globalcontraction}.
\end{proof}

\end{document}
